\begin{frame}{Demostraciones del teorema de compacidad}
   \begin{columns}
      \begin{column}{0.48\textwidth}
         % https://tikzcd.yichuanshen.de/#N4Igdg9gJgpgziAXAbVABwnAlgFyxMJZABgBpiBdUkANwEMAbAVxiRAB12cYAPHYACowIAJxgBbOgAJYUgMYRxaBjDw4mUAL4hNpdJlz5CKMgGYqtRizadufQcLGSZMeYrR05WKHS06LMFAA5vBEoABmIopIZCA4EEgAjJoUmkA
\begin{tikzcd}
\text{Teorema de completitud} \arrow[ddd] \\
                                          \\
                                          \\
\text{Teorema de compacidad}             
\end{tikzcd}

      \end{column}
      \begin{column}{0.48\textwidth}
         \begin{itemize}
            \item <2-> $\mathcal{E}$ es consistente.
            \item <3-> $\mathcal{E}$ puede ser completado: existe $\mathcal{E}^{\#}$ completo y f.s. tal que $\mathcal{E}\subseteq\mathcal{E}^{\#}$
            \item <4-> $\mathcal{E}^{\#}$ es Gödeliano. Por el teorema de completitud $\mathcal{E}^{\#}$ tiene un modelo y por tanto también $\mathcal{E}$.
         \end{itemize}
      \end{column}
   \end{columns}
\end{frame}
\begin{frame}{Demostraciones del teorema de compacidad}
   \begin{columns}
      \begin{column}{0.48\textwidth}
         \begin{tikzcd}
\text{Demostraciones topológicas} \arrow[ddd] \\
                                          \\
                                          \\
\text{Teorema de compacidad}             
\end{tikzcd}

      \end{column}
      \begin{column}{0.48\textwidth}
         \begin{itemize}
            \item Teorema de Tychonoff
            \item Conjunto de Cantor
         \end{itemize}
      \end{column}
   \end{columns}
\end{frame}
\begin{frame}{Demostraciones del teorema de compacidad}
    \begin{columns}
      \begin{column}{0.48\textwidth}
         \begin{tikzcd}
\text{Demostraciones algebraicas} \arrow[ddd] \\
                                          \\
                                          \\
\text{Teorema de compacidad}             
\end{tikzcd}

      \end{column}
      \begin{column}{0.48\textwidth}
         \begin{itemize}
            \item Álgebra de Lindembaum-Tarski
            \item Filtros
         \end{itemize}
      \end{column}
   \end{columns}
\end{frame}
\begin{frame}{Demostraciones del teorema de compacidad}
    \begin{columns}
      \begin{column}{0.48\textwidth}
         \begin{tikzcd}
\text{Demostraciones híbridas} \arrow[ddd] \\
                                          \\
                                          \\
\text{Teorema de compacidad}             
\end{tikzcd}

      \end{column}
      \begin{column}{0.48\textwidth}
         \begin{itemize}
            \item Ultraproductos y espacios de Stone
            \item Espacios de Stone de álgebras boolenas
         \end{itemize}
      \end{column}
   \end{columns}
\end{frame}


